\chapter{Introduction}
\section{Motivation}
\label{ssec:motivation}
High-fidelity computational fluid dynamics (CFD) has become an invaluable tool in aerodynamic design, but the process of manually designing shapes to iterate upon existing geometries leaves  uncertainty as to whether further performance could be gained. With commitments to reduce aircraft emissions by 50\% by 2050 relative to 2005 levels~\cite{tyler_iata_2014}, this uncertainty must be minimized, as efficiency gains will require substantial innovation in many aspects of aircraft design. To this end, CFD coupled with numerical optimization provides an opportunity to improve upon the aerodynamic design process by automatically making shape modifications with the explicit purpose of working towards a specific design objective—a process known as aerodynamic shape optimization (ASO).
In the context of aircraft design, ASO has the added benefit of changing many geometric design variables at once. \par

In decades past, the available computational power, flow solution methods, and numerical optimization methods used in ASO problems severely limited their scope. Early studies included design variables that could be formulated with extremely low cost, such as a wing's angle of attack, along with airfoil shape functions and twist at a handful of stations over a wingspan~\cite{hicks_wing_1978}. 
Some design problems made use of a target pressure coefficient distribution to guide the automated modification of design variables, while others made use of gradient-based methods to minimize some objective, such as the total drag over the geometry of interest. 
As many as 11 variables would be repeatedly modified until the pressure coefficient at the indicated span-wise sections were within a specified tolerance, or until the drag could no longer be reduced without violating design constraints. 
The challenges encountered by these early studies were numerous, but could be primarily attributed to computational resources restricting the size and complexity of inputted geometries, CFD methods available for ASO problems being limited to solving inviscid models such as the Euler equations, and the high cost of numerical methods required for optimization, such as finite-difference methods being used for computing gradients~\cite{hicks_wing_1978}. 
Hicks and Henne briefly highlighted these computational and optimization challenges, stating that planform variables such as aspect ratio, taper, and sweep could be used if general computational power and gradient-computation cost were improved~\cite{hicks_wing_1978}. \par

Since then, advances by Jameson through the adjoint method~\cite{jameson_aerodynamic_1988, jameson_computational_1989} and improved computing capabilities have made ASO on the order of 1000 design variables possible. 
The adjoint method is particularly notable, as it permits the computation of gradients through the solution of a linear system that is independent of the number of design variables. Instead of conducting expensive finite-difference calculations for each design variable, as done by Hicks and Henne, the adjoint method enables a linear system to be solved for the gradients. The number of linear solution computations in each design iteration totals the number of flow-dependent constraints plus one, to account for the objective function. 
This decouples the gradient computation cost from the number of design variables, permitting the resolution of ASO problems to be greatly increased by enabling shape modifications at significantly more spanwise and chordwise points than in the past.
% REMOVED SECTION ON OPTIMIZATION METHODS
% In contrast to inverse methods, the widespread adoption of objective or cost functions in ASO problems enable the design space to become expansive. Jameson highlighted the advantages of this approach in comparison to using target pressure distributions during his derivation of the adjoint method~\cite{jameson_aerodynamic_1988}. Methods using target pressure distributions assume that the target is the optimal distribution subject to certain constraints. However, a target pressure distribution is manually formulated by a designer, such that the resultant target may vary depending on the experience of the designer and the knowledge of the day. For example, the target pressure distribution for a transonic aircraft may have looked very different prior to the commercial adoption of supercritical airfoils in the 1970s, such that the target pressure distribution would be "optimal" but worse performing compared to a supercritical airfoil. Jameson emphasized that worse still is when using inverse methods, a target pressure distribution may not always be feasible even when the designer believes it to be optimal~\cite{jameson_aerodynamic_1988}. For example, a shape for which stagnation pressure is achieved over its entire surface is not possible. In comparison, the use of an objective function permits much greater flexibility in optimizing a design. An objective function allows one to obtain a minimum within the feasible set, regardless of constraints or knowledge of what the optimum pressure distribution or shape properties may look like. Moreover, given the flexibility in its constraints, the minimization of an objective function could explore nontraditional shapes while searching for an optimum. In this way, the widespread adoption of objective functions in ASO greatly enabled the optimization of wings in an investigative manner, where details of an optimal result may not be known beforehand~\cite{jameson_optimum_1995}, thereby enabling designers to minimize any single or combination of performance metrics without the risk of wasted resources if the desired pressure distribution was not feasible ~\cite{martins_aerodynamic_2022}.

The expansion of ASO methods to include viscous and turbulent flow solutions advanced design capabilities further, where the solution of the Reynolds-Averaged Navier Stokes coupled with the Spalart-Allmaras turbulence model enables viscous sources of drag to be accounted for and reduced during optimization~\cite{jameson_optimum_1998,nielsen_aerodynamic_1999}. 
This was a notable development, as critical features such as flow separation are not modelled by an inviscid flow solution. An inviscid ASO result could produce a highly cambered trailing edge or high angle of attack, whereas a viscous ASO solution would model the separation and drag that could result from such configurations, steering the wing towards better performing geometries.
In parallel with progress in flow, gradient, and optimization advancements, geometry parameterization, control, and mesh deformation have experienced significant developments as well.
The ability to conduct efficient and precise changes to wing configurations during optimization has been enabled by highly advanced methods such as free-form and axial deformation (FFD)~\cite{gagnon_two-level_2015}, domain-element methods~\cite{morris_domain-element_2009}, B-spline geometry parameterization~\cite{yu_application_1995}, and more.
Numerous combinations of these methods have enabled modern ASO to be conducted with increasingly diverse design variables at incredibly fine detail.
Whereas early studies were limited to optimizing airfoil section shapes through parameterizing functions, modern optimization can be conducted on the section shapes, angle of attack, twist, taper, sweep, and dihedral of a wing all at once, with analysis at tens or even hundreds of chordwise and spanwise points along a wing~\cite{lyu_benchmarking_2014, zingg_aerodynamic_2017, streuber_evaluating_2021, reist_aerodynamic_2013, reist_optimization_2015, reist_aerodynamic_2016, kuntawala_preliminary_2011}. \par

Indeed, today's optimizations are capable of conducting investigations of configurations in viscous flow with extremely high numbers of degrees of freedom, which is a stark improvement over the optimizations of days past.
Yet, despite these advancements, ASO that includes wing sweep as a design variable is rare.
Wing sweep is an important design variable through its influence on wave drag, where an increased wing sweep could reduce or eliminate wave drag, and vice versa.
This is especially relevant to medium and long-haul transport aircraft, as they spend a significant portion of time in transonic cruise conditions during which wave drag may occur.
Yet, many studies of conventional transport aircraft wing configurations either do not include sweep at all~\cite{jameson_optimum_1995, jameson_optimum_1998, nielsen_aerodynamic_1999, reuther_aerodynamic_1996, kenway_cad-free_2010, ledoux_study_2015, kenway_multipoint_2015, martins_high-fidelity_2017, chen_aerodynamic_2015, telidetzki_application_2014, lyu_rans-based_2014,  lyu_aerodynamic_2015, osusky_drag_2015, chen_aerodynamic_2016, meheut_gradient-based_2016, shi-dong_adjoint-based_2017, shi-dong_approximate_2018, yu_influence_2018, mangano_multipoint_2019, poole_objectives_2020} or study them only at a limited range of fixed, discrete values~\cite{darshan_aerodynamic_2021}.
A rare example of pure ASO on a conventional configuration with wing sweep analyzed on a continuous basis is detailed in Koo and Zingg~\cite{koo_investigation_2018}, who made use of free-form deformation (FFD) volumes and axial curves to conduct a transonic viscous ASO of various wings. However, the optimized wing sweep by Koo and Zingg was driven to its upper bound on multiple planform geometries, and it is unclear how the wing sweep and other design variables would behave given an expanded sweep range or an inactive sweep constraint. Moreover, the relationship between aircraft drag and wing sweep angle is poorly characterized when other design variables are included. Should an optimal wing sweep angle exist, it is unclear whether tangible performance improvements can be continuously obtained as one approaches the optimum. Darshan et al. conducted ASO at discrete wing sweep angles on the NASA CRM wing, but only considered angles below the original sweep angle of 37.5\degree\ ~\cite{darshan_aerodynamic_2021}. Therefore, it is unclear what the relationship between drag and sweep angle is at greater sweep angles, and how this relationship may change as the sweep angle increases.\par

It is important to note that a great number of ASO studies have been conducted with many other design variables excluding wing sweep. A thorough search of relevant literate did not reveal any explicit justification for excluding wing sweep as a design variable. It is possible that the exclusion of wing sweep is due to the effects that wing sweep modifications have on non-aerodynamic factors in aircraft design.
Changes in an aircraft's wing sweep angle will result in changes to its overall weight and stability characteristics~\cite{raymer_aircraft_2018}.
Hence, a purely aerodynamic optimization of a wing sweep angle is not realistic in the context of an aircraft's overall performance in multiple disciplines, which may dissuade some from excluding it as a design variable.\par

In looking towards modern multi-disciplinary aircraft optimization studies, one finds that wing sweep is occasionally included~\cite{jameson_aerodynamic_2003, kenway_cad-free_2010, jameson_multi-point_2007}, but increases in sweep are penalized by the corresponding increases in structural strength requirements and weight, imposing a direct influence on the resultant sweep angle outside of aerodynamic effects.
Kenway et al.~\cite{kenway_cad-free_2010} noted this in their MDO study of the DPW4 wing-body-tail configuration, where the optimized sweep showed a trade-off between structural and aerodynamic requirements. It was noted that "from a strictly aerodynamic perspective the wing should sweep to its aft limit of 45 degrees"~\cite{kenway_cad-free_2010}, but further insight or evidence supporting this comment are not provided. Furthermore, no comments were made on such behaviour if the sweep limit were expanded further. \par

The uncertainty of the behaviour of wing sweep in a pure ASO is further compounded by the rarity of wing sweep in studies of ASO multimodality. Numerous studies have analyzed the impact of additional degrees of freedom in the design space of aircraft ASO problems, where planform, sectional, and twist variables have been added in succession and optimized—though few consider sweep as a major focus~\cite{lyu_rans-based_2014,  osusky_drag_2015, chernukhin_multimodality_2013, streuber_investigation_2017, koo_investigation_2018, poole_global_2018, poole_identifying_2018, streuber_evaluating_2021}. These studies have concluded that increased dimensionality in ASO generally results in multimodality and the generation of local optima regardless of constraints~\cite{lyu_rans-based_2014,  osusky_drag_2015, chernukhin_multimodality_2013, streuber_investigation_2017, koo_investigation_2018, poole_global_2018, poole_identifying_2018, streuber_evaluating_2021}. Studies that include wing sweep study it in cases that are less extensive; nonetheless, it is seen that the addition of wing sweep may impose greater difficulty locating a high-quality optimum due to the high likelihood of multimodality that results from its addition as a design variable~\cite{streuber_investigation_2017, poole_global_2018, poole_identifying_2018, streuber_evaluating_2021}. However, further examinations on the performance of optima as they pertain to specific sweep angles were not conducted, and it is unclear whether the sweep angle constraints are active at the resultant optima. A thorough search of the relevant literature revealed no further information in ASO, multimodality, or MDO studies on how the wing sweep angle in a conventional wing configuration would behave given a very high upper bound.

In addition to its importance in the design of conventional commercial aircraft configurations, optimization of the wing sweep angle is of great importance in the innovative design of sustainable aircraft. New developments in wing design such as the use of natural laminar flow (NLF) have the potential to substantially reduce the friction drag on an aircraft by maintaining a laminar boundary layer~\cite{tokugawa_natural_2023}. However, NLF design is significantly more difficult when applied to swept wings. A swept wing induces a crossflow, which typically dominates as the source of boundary layer instability in swept wing NLF design~\cite{tokugawa_natural_2023}. In contrast, Tollmien-Schlichting instability dominates in straight wings, as the wing does not induce a crossflow~\cite{tokugawa_natural_2023}. Sources of boundary layer instability contribute to boundary layer transition from laminar to turbulent flow, the latter of which is undesirable in a NLF wing~\cite{tokugawa_natural_2023}. Studies of NLF wing design have emphasized the importance of including wing sweep as a design variable so as to maximize the extent of NLF obtained~\cite{tokugawa_natural_2023,campbell_natural_2016,lynde_computational_2017,sudhi_design_2023}. Methods exist by which the laminar boundary layer can be maintained and instability reduced at higher sweep angles, but these methods require active or hybrid flow control~\cite{campbell_natural_2016, lynde_computational_2017}, which are associated with numerous challenges beyond laminar flow design that are out of the scope of this work.

Lastly, innovations in other areas such as blended-wing-body (BWB) aircraft design have demonstrated the viability and importance of including wing sweep angles among wide ranges of design variables~\cite{kuntawala_preliminary_2011,reist_aerodynamic_2013,reist_optimization_2015,reist_aerodynamic_2016}. Studies such as that by Kuntawala et al.~\cite{kuntawala_preliminary_2011} have conducted pure ASO on BWBs in which drag reductions of over 5\% were observed between geometries where the only major differences occurred in wing twist and sweep angles. However, even BWB studies with wing sweep have displayed instances of the optimized sweep angle being driven to the maximum extent, such as that conducted by Reist and Zingg on a BWB regional transport aircraft optimization~\cite{reist_aerodynamic_2013}. Nevertheless, these studies demonstrate that the use of a wide range of design variables has the potential to produce significant reductions in drag, and that the wing sweep angle may be an important factor contributing to this. The same may be true for conventional configurations, but a search of relevant literature yielded no further studies with a similarly extensive range of design variables.

Therefore, it is of great interest to observe the resultant wing sweep angle of a pure ASO study on conventional aircraft configurations. A study of this nature would act as a demonstration of modern optimization methodology in its ability to reliably optimize wing sweep. Furthermore, the resultant configuration would provide a means for analyzing the changes in total drag or components such as friction or pressure drag once wave drag is substantially reduced or negligible as a result of wing sweep-inclusive ASO. Lastly, a conventional configuration resulting from pure ASO with wing sweep will serve as a point of comparison between optimized designs for conventional configurations, NLF configurations, and BWB configurations with respect to drag minimization from a purely aerodynamic perspective. In particular, since a configuration designed for NLF must include the wing sweep angle as a design variable, an optimized conventional aircraft with wing sweep also included as a design variable would enable a more direct performance comparison between the two.

\section{Objectives}
The objective of this study is to build upon the findings of Koo and Zingg~\cite{koo_investigation_2018} and Streuber and Zingg~\cite{streuber_investigation_2017, streuber_evaluating_2021} in two areas, First, the behaviour of the wing sweep angle in the aerodynamic shape optimization of multiple wing configurations when given very high upper bounds will be investigated. This study will focus on the NASA Common Research Model (CRM) and Airbus A320neo planforms, enabling the analysis of optimization behaviour with wing sweep when advanced configurations are used as initial geometries. This will be achieved by characterizing the convergence behaviour of the ASO on the aforementioned wing configurations when wing sweep is included as design variable, as well as through performance analysis of the resulting optimum designs. The ASOs performed will include typical design variables such as wing angle of attack, section shape at multiple points along the wingspan, twist, and more. A second study will conduct ASO on the NASA CRM at fixed sweep angles. This study will enable the characterization of the relationship between drag and sweep angles when multiple design variables are included. The ASO performed will include include the same design variables as the first study.